% !TeX root = ../main.tex
\section{Related Work}
The modeling and identification of flapping-wing micro air vehicles (FWMAVs) have been studied for controller design, simulation, and flight performance analysis. A first group of works follows the classical aircraft system-identification approach and adapts it to the periodic dynamics of flapping flight. For the DelFly platform, Caetano et al. reconstructed flap-averaged aerodynamic forces and moments from flight-test data and identified linear aerodynamic models~\cite{caetano_linear_2013}. Armanini et al. identified black-box LTI and time-varying models from free-flight data, showing the relevance of periodic effects in the vehicle dynamics~\cite{armanini_black-box_2015,armanini_time-varying_2016}. Grey-box models and onboard-sensor-based methods have also been reported for tailless platforms and longitudinal motion~\cite{nijboer_longitudinal_2020,aurecianus_longitudinal_2022}. These methods are useful because they produce compact and interpretable models. However, they usually rely on linearized, flap-averaged, or low-dimensional descriptions of the aerodynamic loads.

Data-driven aerodynamic models provide another way to describe the unsteady forces produced by flapping wings. Artificial neural networks were first used to predict lift and drag histories for prescribed flapping airfoil motions~\cite{kurtulus_ability_2009}. Later works used neural networks, Gaussian processes, and reduced-order models to accelerate the evaluation or optimization of hovering and prescribed-kinematics flapping wings~\cite{nguyen_neural-network-based_2019,calado_robust_2023,lan_neural_2022}. More recently, deep surrogate models have been used for time-history performance prediction, active-learning-based kinematic optimization, and experimental lift prediction of flapping-wing rotors~\cite{wang_time-history_2024,corban_discovering_2023,chen_data-driven_2026}. These studies show that data-driven models can represent nonlinear and history-dependent aerodynamic effects. However, the training data are commonly generated from CFD, DNS, wind-tunnel measurements, prescribed wing motions, or bench-top experiments. The resulting models are therefore not directly built from onboard free-flight logs of a complete vehicle.

Learning has also been introduced in FWMAV dynamics, estimation, and control. Adaptive neural networks and RBFNNs have been used to compensate unknown aerodynamic perturbations and nonlinear coupling in trajectory tracking and position control~\cite{he_adaptive_2017,he_modeling_2021,qian_neural_2023}. Gaussian-process models have been used to learn flapping-wing robot dynamics during trajectory tracking, and to estimate periodic components in state-estimation pipelines~\cite{qian_learning_2023,qian_learning_2026}. Simulation tools have further combined system identification, flapping-wing dynamics, and reinforcement-learning control for hummingbird-scale robots~\cite{fei_flappy_2019}. These works indicate that learning methods are effective for uncertainty compensation and controller synthesis. Nevertheless, the learned model is often used as part of a controller, estimator, or simulator, rather than as a flight-log-based predictor of the vehicle aerodynamic wrench.

In this letter, we study a complementary problem. Instead of constructing an aerodynamic surrogate from prescribed wing kinematics, or using a neural network only for control compensation, we identify a mapping from onboard flight logs to the body-frame effective aerodynamic wrench. The inputs include the vehicle state, actuator commands, and flapping phase and frequency information. The output is treated as a time-varying force-and-moment signal, so that the model can account for the coupled effects of body motion, tail actuation, wing phase, and flapping-rate variation. This formulation allows different feed-forward and sequence backbones, including MLP, GRU, TCN, and Transformer models, to be evaluated under the same data pipeline. Thus, the proposed study complements low-order flight-dynamics identification and controlled aerodynamic surrogate modeling by focusing on sequence modeling of aerodynamic forces and moments from real free-flight data of a bird-scale flapping-wing UAV.
